% Chunk: physical PDF pages 475--479 (printed pages 452--456).
% Semantic coverage: conclusion of Section 10.5; Section 10.6 through
% Eq. (10.6.23), with its final continuation on physical page 480 excluded.
% Equations: (10.5.19)--(10.5.21), (10.6.1)--(10.6.23).
% Figures: none. Footnotes: one. Uncertainties: none.

still at $q^2=0$, indicating that \emph{radiative corrections do not give the
photon a mass}.

For a renormalized electromagnetic field, radiative corrections should also
not alter the gauge-invariant part of the residue of the photon pole in
Eq.~\eqref{eq:10.5.17}, so
\begin{equation}
  \pi(0)=0.
  \tag{10.5.19}\label{eq:10.5.19}
\end{equation}
This condition leads to a determination of the electromagnetic field
renormalization constant $Z_3$. Recall that when expressed in terms of the
renormalized field~\eqref{eq:10.4.17}, the electrodynamic Lagrangian takes the
form
\begin{equation*}
  \mathcal L=-\frac14 Z_3(\partial_\mu A_\nu-\partial_\nu A_\mu)
  (\partial^\mu A^\nu-\partial^\nu A^\mu)
  +\mathcal L_M\bigl(\Psi_\ell,
  [\partial_\mu-iZ_3q_\ell A_\mu]\Psi_\ell\bigr).
\end{equation*}
The function $\pi(q^2)$ in the one-photon-irreducible amplitude is then
\begin{equation}
  \pi(q^2)=1-Z_3+\pi_{\mathrm{LOOP}}(q^2),
  \tag{10.5.20}\label{eq:10.5.20}
\end{equation}
where $\pi_{\mathrm{LOOP}}$ is the contribution of loop diagrams. It follows
that
\begin{equation}
  Z_3=1+\pi_{\mathrm{LOOP}}(0).
  \tag{10.5.21}\label{eq:10.5.21}
\end{equation}
In practice, we just calculate the loop contributions and subtract a constant
in order to make $\pi(0)$ vanish.

Incidentally, Eq.~\eqref{eq:10.5.18} shows that for $q^2\neq0$ the gauge term
in the photon propagator is altered by radiative corrections. The one
exception is the case of Landau gauge, for which $\bar\xi=\xi=1$ for all
$q^2$.

\subsection{Electromagnetic Form Factors and Magnetic Moment}

Suppose that we want to calculate the scattering of a particle by an external
electromagnetic field (or by the electromagnetic field of another particle),
to first order in this electromagnetic field, but to all orders in all other
interactions (including electromagnetic) of our particle. For this purpose,
we need to know the sum of the contributions of all Feynman diagrams with one
incoming and one outgoing particle line, both on the mass shell, plus a photon
line, which may be on or off the mass shell. According to the theorem of
Section 6.4, this sum is given by the one-particle matrix element of the
electromagnetic current $J^\mu(x)$. Let us see what governs the general form
of this matrix element.

According to spacetime translation invariance, the one-particle matrix element
of the electromagnetic current takes the form
\begin{equation}
  \bra{p',\sigma'}J^\mu(x)\ket{p,\sigma}
  =\exp\bigl(i(p-p')\cdot x\bigr)
  \bra{p',\sigma'}J^\mu(0)\ket{p,\sigma}.
  \tag{10.6.1}\label{eq:10.6.1}
\end{equation}
The current conservation condition $\partial_\mu J^\mu=0$ then requires
\begin{equation}
  (p'-p)_\mu\bra{p',\sigma'}J^\mu(0)\ket{p,\sigma}=0.
  \tag{10.6.2}\label{eq:10.6.2}
\end{equation}

Also, setting $\mu=0$ and integrating over all $\mathbf{x}$ gives
\begin{equation*}
  \bra{p',\sigma'}Q\ket{p,\sigma}
  =(2\pi)^3\delta^3(\mathbf{p}-\mathbf{p}')
  \bra{p',\sigma'}J^0(0)\ket{p,\sigma}.
\end{equation*}
Using Eq.~\eqref{eq:10.4.8}, this gives
\begin{equation}
  \bra{p',\sigma'}J^0(0)\ket{p,\sigma}
  =(2\pi)^{-3}q\,\delta_{\sigma'\sigma},
  \tag{10.6.3}\label{eq:10.6.3}
\end{equation}
where $q$ is the particle charge.

We also have at our disposal the constraints on the current matrix elements
imposed by Lorentz invariance. To explore these, we will limit ourselves to
the simplest cases: spin zero and spin $\frac12$. The analysis presented here
provides an example of techniques that are useful for other currents, such as
those of the semi-leptonic weak interactions.

\begin{center}
  \emph{Spin Zero}
\end{center}

For spin zero, Lorentz invariance requires the one-particle matrix element of
the current to take the general form
\begin{equation}
  \bra{p'}J^\mu(0)\ket{p}
  =q(2\pi)^{-3}(2p'^0)^{-1/2}(2p^0)^{-1/2}\mathcal J^\mu(p',p),
  \tag{10.6.4}\label{eq:10.6.4}
\end{equation}
where $p^0$ and $p'^0$ are the mass-shell energies
($p^0=\sqrt{\mathbf{p}^2+m^2}$), and $\mathcal J^\mu(p',p)$ is a four-vector
function of the two four-vectors $p'^\mu$ and $p^\mu$. (We have extracted a
factor of the charge $q$ of the particle from $\mathcal J$ for future
convenience.) Obviously, the most general such four-vector function takes the
form of a linear combination of $p'^\mu$ and $p^\mu$, or equivalently of
$p'^\mu+p^\mu$ and $p'^\mu-p^\mu$, with scalar coefficients. But the scalars
$p^2$ and $p'^2$ are fixed at the values $p'^2=p^2=-m^2$, so the scalar
variables that can be formed from $p^\mu$ and $p'^\mu$ can be taken as
functions only of $p\cdot p'$, or equivalently of
\begin{equation}
  k^2\equiv(p-p')^2=-2m^2-2p\cdot p'.
  \tag{10.6.5}\label{eq:10.6.5}
\end{equation}
Thus the function $\mathcal J^\mu(p',p)$ must take the form
\begin{equation}
  \mathcal J^\mu(p',p)
  =(p'+p)^\mu F(k^2)+i(p'-p)^\mu H(k^2).
  \tag{10.6.6}\label{eq:10.6.6}
\end{equation}
The fact that $J^\mu$ is Hermitian implies that
$\mathcal J^\mu(p',p)^*=\mathcal J^\mu(p,p')$, so that both $F(k^2)$ and
$H(k^2)$ are real.

Now $(p'-p)\cdot(p'+p)$ vanishes, while $(p'-p)^2=k^2$ is not generally
zero, so the condition of current conservation is simply
\begin{equation}
  H(k^2)=0.
  \tag{10.6.7}\label{eq:10.6.7}
\end{equation}
Also, setting $\mathbf{p}'=\mathbf{p}$ and $\mu=0$ in
Eq.~\eqref{eq:10.6.4}, and comparing with Eq.~\eqref{eq:10.6.3}, we find that
\begin{equation}
  F(0)=1.
  \tag{10.6.8}\label{eq:10.6.8}
\end{equation}
The function $F(k^2)$ is called the \emph{electromagnetic form factor} of the
particle.

\begin{center}
  \emph{Spin $\frac12$}
\end{center}

For spin $\frac12$, Lorentz invariance requires the one-particle matrix element
of the current to take the general form
\begin{equation}
  \bra{p',\sigma'}J^\mu(0)\ket{p,\sigma}
  =iq(2\pi)^{-3}\bar u(p',\sigma')\Gamma^\mu(p',p)u(p,\sigma)
  \tag{10.6.9}\label{eq:10.6.9}
\end{equation}
where $\Gamma^\mu$ is a four-vector $4\cdot4$ matrix function of $p'^\nu$,
$p^\nu$, and $\gamma^\nu$, and $u$ is the usual Dirac coefficient function.
We have extracted a factor $iq$ to make the normalization of $\Gamma^\mu$ the
same as in the previous section.

Just as for any $4\cdot4$ matrix, we may expand $\Gamma^\mu$ in the 16
covariant matrices $\mathbf{1}$, $\gamma_\rho$,
$[\gamma_\rho,\gamma_\sigma]$, $\gamma_5\gamma_\rho$, and $\gamma_5$. The most
general four-vector $\Gamma^\mu$ can therefore be written as a linear
combination of
\begin{align*}
  \mathbf{1}:\quad& p^\mu,\quad p'^\mu,
  \\
  \gamma_\rho:\quad& \gamma^\mu,\quad
    p^\mu\sl{p},\quad p'^\mu\sl{p},\quad
    p^\mu\sl{p}',\quad p'^\mu\sl{p}',
  \\
  [\gamma_\rho,\gamma_\sigma]:\quad&
    [\gamma^\mu,\sl{p}],\quad [\gamma^\mu,\sl{p}'],\quad
    [\sl{p},\sl{p}']p^\mu,\quad [\sl{p},\sl{p}']p'^\mu,
  \\
  \gamma_5\gamma_\rho:\quad&
    \gamma_5\gamma_\rho\epsilon^{\rho\mu\nu\sigma}p_\nu p'_\sigma,
  \\
  \gamma_5:\quad& \text{NONE}
\end{align*}
with the coefficient of each term a function of the only scalar variable in
the problem, the quantity~\eqref{eq:10.6.5}. This can be greatly simplified by
using the Dirac equations satisfied by $u$ and $\bar u$:
\begin{equation*}
  \bar u(p',\sigma')(i\sl{p}'+m)=0,
  \qquad
  (i\sl{p}+m)u(p,\sigma)=0.
\end{equation*}
In consequence, we can drop\footnote{This is obvious for the terms
$p^\mu\sl{p}$, $p'^\mu\sl{p}$, $p^\mu\sl{p}'$, and
$p'^\mu\sl{p}'$, which may be replaced respectively with $imp^\mu$,
$imp'^\mu$, $imp^\mu$, and $imp'^\mu$, which are the same as terms already
on our list. Also, we can write
\[
  [\gamma^\mu,\sl{p}]
  =2\gamma^\mu\sl{p}-\{\gamma^\mu,\sl{p}\}
  =2\gamma^\mu\sl{p}-2p^\mu,
\]
which may be replaced with $2im\gamma^\mu-2p^\mu$, a linear combination of
terms already on our list. The same applies to
$[\gamma^\mu,\sl{p}']$. Also,
\[
  [\sl{p},\sl{p}']
  =-2\sl{p}'\sl{p}+\{\sl{p},\sl{p}'\}
  =-2\sl{p}'\sl{p}+2p\cdot p',
\]
which may be replaced with $2m^2+2p\cdot p'=-k^2$. Hence the terms
$[\sl{p},\sl{p}']p^\mu$ and $[\sl{p},\sl{p}']p'^\mu$ give nothing new.
Finally, to deal with the last term we may use the relation
\[
  \gamma_5\gamma_\rho\epsilon^{\rho\mu\nu\sigma}
  =\frac{i}{6}\bigl(
  \gamma^\mu\gamma^\nu\gamma^\sigma
  +\gamma^\sigma\gamma^\mu\gamma^\nu
  +\gamma^\nu\gamma^\sigma\gamma^\mu
  -\gamma^\mu\gamma^\sigma\gamma^\nu
  -\gamma^\nu\gamma^\mu\gamma^\sigma
  -\gamma^\sigma\gamma^\nu\gamma^\mu\bigr).
\]
Contracting this with $p_\nu$ and $p'_\sigma$ and then moving all $\sl{p}$
factors to the right and $\sl{p}'$ factors to the left again gives a linear
combination of $p^\mu$, $p'^\mu$, and $\gamma^\mu$.} all but the first three
entries: $p^\mu$, $p'^\mu$, and $\gamma^\mu$. We conclude that, on the fermion
mass shell, $\Gamma^\mu$ may be expressed as a linear combination of
$\gamma^\mu$, $p^\mu$, and $p'^\mu$, which we choose to write as
\begin{align}
  \bar u(p',\sigma')\Gamma^\mu(p',p)u(p,\sigma)
  =\bar u(p',\sigma')\biggl[
  &\gamma^\mu F(k^2)-\frac{i}{2m}(p+p')^\mu G(k^2) \notag\\
  &+\frac{(p-p')^\mu}{2m}H(k^2)\biggr]u(p,\sigma).
  \tag{10.6.10}\label{eq:10.6.10}
\end{align}

The hermiticity of $J^\mu(0)$ implies that
\begin{equation}
  \beta\Gamma^{\mu\dagger}(p',p)\beta=-\Gamma^\mu(p,p'),
  \tag{10.6.11}\label{eq:10.6.11}
\end{equation}
so that $F(k^2)$, $G(k^2)$, and $H(k^2)$ must all be \emph{real} functions of
$k^2$.

The conservation condition~\eqref{eq:10.6.2} is automatically satisfied by
the first two terms in Eq.~\eqref{eq:10.6.10}, because
\begin{equation*}
  (p'-p)_\mu\gamma^\mu
  =-i\bigl[(i\sl{p}'+m)-(i\sl{p}+m)\bigr]
\end{equation*}
and
\begin{equation*}
  (p'-p)\cdot(p'+p)=p'^2-p^2.
\end{equation*}
On the other hand, $(p'-p)^2$ does not in general vanish, so current
conservation requires the third term to vanish
\begin{equation}
  H(k^2)=0.
  \tag{10.6.12}\label{eq:10.6.12}
\end{equation}
Also, letting $\mathbf{p}'\to\mathbf{p}$ in Eqs.~\eqref{eq:10.6.9} and
\eqref{eq:10.6.10}, we find
\begin{equation*}
  \bra{p,\sigma'}J^\mu(0)\ket{p,\sigma}
  =iq(2\pi)^{-3}\bar u(p,\sigma')
  \left[\gamma^\mu F(0)-\frac{i}{m}p^\mu G(0)\right]u(p,\sigma).
\end{equation*}
Using the identity $\{\gamma^\mu,i\sl{p}+m\}=2m\gamma^\mu+2ip^\mu$,
we also have
\begin{equation*}
  \bar u(p,\sigma')\gamma^\mu u(p,\sigma)
  =-\frac{ip^\mu}{m}\bar u(p,\sigma')u(p,\sigma).
\end{equation*}
Recall also that
\begin{equation*}
  \bar u(p,\sigma')u(p,\sigma)=\delta_{\sigma'\sigma}m/p^0
\end{equation*}
and therefore
\begin{equation}
  \bra{p,\sigma'}J^\mu(0)\ket{p,\sigma}
  =q(2\pi)^{-3}(p^\mu/p^0)\delta_{\sigma'\sigma}[F(0)+G(0)].
  \tag{10.6.13}\label{eq:10.6.13}
\end{equation}
Comparing this with Eq.~\eqref{eq:10.6.3} yields the normalization condition
\begin{equation}
  F(0)+G(0)=1.
  \tag{10.6.14}\label{eq:10.6.14}
\end{equation}

It may be useful to note that the electromagnetic vertex matrix $\Gamma^\mu$
is commonly written in terms of two other matrices, as
\begin{align}
  \bar u(p',\sigma')\Gamma^\mu(p',p)u(p,\sigma)
  =\bar u(p',\sigma')\biggl[&\gamma^\mu F_1(k^2) \notag\\
  &+\frac{i}{2}[\gamma^\mu,\gamma^\nu](p'-p)_\nu F_2(k^2)
  \biggr]u(p,\sigma).
  \tag{10.6.15}\label{eq:10.6.15}
\end{align}
As already mentioned, we may rewrite the matrix appearing in the second term
in terms of those used in defining $F(k^2)$ and $G(k^2)$:
\begin{align}
  &\bar u(p',\sigma')\frac{i}{2}[\gamma^\mu,\gamma^\nu]
  (p'-p)_\nu u(p,\sigma) \notag\\
  &\quad=\bar u(p',\sigma')\bigl[
  -i\sl{p}'\gamma^\mu+\tfrac{i}{2}\{\gamma^\mu,\sl{p}'\}
  -i\gamma^\mu\sl{p}+\tfrac{i}{2}\{\gamma^\mu,\sl{p}\}
  \bigr]u(p,\sigma) \notag\\
  &\quad=\bar u(p',\sigma')
  [i(p^\mu+p'^\mu)+2m\gamma^\mu]u(p,\sigma).
  \tag{10.6.16}\label{eq:10.6.16}
\end{align}

Comparing Eq.~\eqref{eq:10.6.15} with Eq.~\eqref{eq:10.6.10}, we find
\begin{equation}
  F(k^2)=F_1(k^2)+2mF_2(k^2)
  \tag{10.6.17}\label{eq:10.6.17}
\end{equation}
\begin{equation}
  G(k^2)=-2mF_2(k^2).
  \tag{10.6.18}\label{eq:10.6.18}
\end{equation}
The normalization condition~\eqref{eq:10.6.14} now reads
\begin{equation*}
  F_1(0)=1.
\end{equation*}

In order to evaluate the magnetic moment of our particle in terms of its form
factors, let us consider the spatial part of the vertex function in the case
of small momenta, $|\mathbf{p}|,|\mathbf{p}'|\ll m$. For this purpose, it is
useful to use Eq.~\eqref{eq:10.6.16} to rewrite Eq.~\eqref{eq:10.6.10} (with
$H=0$) in a third form:
\begin{align}
  \bar u(p',\sigma')\Gamma^\mu(p',p)u(p,\sigma)
  =-\frac{i}{2m}\bar u(p',\sigma')\biggl[&
  (p+p')^\mu\{F(k^2)+G(k^2)\} \notag\\
  &-\frac12[\gamma^\mu,\gamma^\nu](p'-p)_\nu F(k^2)
  \biggr]u(p,\sigma).
  \tag{10.6.19}\label{eq:10.6.19}
\end{align}
For zero momenta the matrix elements of the commutators of Dirac matrices are
given by Eqs.~\eqref{eq:5.4.19} and~\eqref{eq:5.4.20} as
\begin{equation*}
  \bar u(0,\sigma')[\gamma^i,\gamma^j]u(0,\sigma)
  =4i\epsilon_{ijk}(J_k^{(1/2)})_{\sigma'\sigma},
  \qquad
  \bar u(0,\sigma')[\gamma^i,\gamma^0]u(0,\sigma)=0,
\end{equation*}
where $\mathbf{J}^{(1/2)}=\frac12\boldsymbol{\sigma}$ is the angular momentum
matrix for spin $\frac12$. Hence to first order in the small momenta,
\begin{equation}
  \bar u(p',\sigma')\Gamma^i(p',p)u(p,\sigma)
  \longrightarrow\frac1m
  [ (\mathbf{p}-\mathbf{p}')\cdot\mathbf{J}^{(1/2)}]_{i\sigma'\sigma}F(0).
  \tag{10.6.20}\label{eq:10.6.20}
\end{equation}
In a very weak time-independent external vector potential $\mathbf{A}(\mathbf{x})$
the matrix element of the interaction Hamiltonian
$H'=-\int d^3x\,\mathbf{J}(\mathbf{x})\cdot\mathbf{A}(\mathbf{x})$ between
one-particle states of small momentum is therefore
\begin{align}
  \bra{p',\sigma'}H'\ket{p,\sigma}
  ={}&-\frac{iqF(0)}{m(2\pi)^3}\int d^3x\,
  e^{i(\mathbf{p}-\mathbf{p}')\cdot\mathbf{x}}
  \mathbf{A}(\mathbf{x})\cdot
  [ (\mathbf{p}-\mathbf{p}')\cdot\mathbf{J}^{(1/2)}]_{\sigma'\sigma}
  \notag\\
  ={}&-\frac{qF(0)}{m(2\pi)^3}\int d^3x\,
  e^{i(\mathbf{p}-\mathbf{p}')\cdot\mathbf{x}}
  (\mathbf{J}^{(1/2)})_{\sigma'\sigma}\cdot\mathbf{B}(\mathbf{x}),
  \tag{10.6.21}\label{eq:10.6.21}
\end{align}
where $\mathbf{B}=\boldsymbol{\nabla}\cdot\mathbf{A}$ is the magnetic field.
Hence in the limit of a slowly varying weak magnetic field, the matrix element
of the interaction Hamiltonian is
\begin{equation}
  \bra{p',\sigma'}H'\ket{p,\sigma}
  =-\frac{qF(0)}m(\mathbf{J}^{(1/2)})_{\sigma'\sigma}\cdot\mathbf{B}\,
  \delta^3(\mathbf{p}-\mathbf{p}').
  \tag{10.6.22}\label{eq:10.6.22}
\end{equation}
The magnetic moment $\mu$ for an arbitrary particle of general spin $j$ is
defined by the statement that the matrix element of the interaction of the
particle with a weak static slowly varying magnetic field is
\begin{equation}
  \bra{p',\sigma'}H'\ket{p,\sigma}
  =-\frac{\mu}{j}(\mathbf{J}^{(j)})_{\sigma'\sigma}\cdot\mathbf{B}\,
  \delta^3(\mathbf{p}-\mathbf{p}').
  \tag{10.6.23}\label{eq:10.6.23}
\end{equation}
